\chapter*{Bibliographie et webographie}
\addcontentsline{toc}{chapter}{Bibliographie et webographie}

\begin{thebibliography}{99}

\bibitem{dragon}
A.~V.~Aho, M.~S.~Lam, R.~Sethi, J.~D.~Ullman,
\emph{Compilers: Principles, Techniques, and Tools} (le « Dragon Book »),
2\textsuperscript{e}~édition, Pearson/Addison-Wesley, 2006.

\bibitem{levine}
J.~Levine,
\emph{flex \& bison: Text Processing Tools},
O'Reilly Media, 2009.

\bibitem{flex}
\emph{Flex : The Fast Lexical Analyzer Generator}, manuel de référence.
\url{https://github.com/westes/flex} (consulté en 2026).

\bibitem{bison}
\emph{GNU Bison Manual}, Free Software Foundation.
\url{https://www.gnu.org/software/bison/manual/} (consulté en 2026).

\bibitem{knuth}
D.~E.~Knuth,
\emph{On the Translation of Languages from Left to Right},
Information and Control, vol.~8, n\textsuperscript{o}~6, p.~607--639, 1965.

\bibitem{chomsky}
N.~Chomsky,
\emph{Three Models for the Description of Language},
IRE Transactions on Information Theory, vol.~2, n\textsuperscript{o}~3, 1956.

\bibitem{lesk}
M.~E.~Lesk, E.~Schmidt,
\emph{Lex : A Lexical Analyzer Generator},
Bell Laboratories, Computing Science Technical Report n\textsuperscript{o}~39, 1975.

\bibitem{johnson}
S.~C.~Johnson,
\emph{Yacc: Yet Another Compiler-Compiler},
Bell Laboratories, Computing Science Technical Report n\textsuperscript{o}~32, 1975.

\bibitem{isoc}
ISO/IEC 9899:2011, \emph{Programming languages : C} (norme C11).

\bibitem{intel}
Intel Corporation,
\emph{Intel 64 and IA-32 Architectures Software Developer's Manual},
volume~2 (Instruction Set Reference).

\bibitem{sysv}
\emph{System V Application Binary Interface : Intel386 Architecture Processor
Supplement} (convention d'appel cdecl).

\bibitem{cours}
P.~Kislin-Duval,
\emph{Interprétation \& Compilation}, support de cours, Licence Informatique,
IED, Université Paris~8.

\end{thebibliography}
